\documentclass[11pt,a4paper]{article}
\usepackage[T1]{fontenc}
\usepackage[utf8]{inputenc}
\usepackage[ngerman]{babel}
\usepackage[margin=2.5cm]{geometry}
\usepackage[hidelinks]{hyperref}
\setlength{\parindent}{0pt}
\setlength{\emergencystretch}{2em}
\begin{document}
\begin{center}
  {\Large\textbf{Webtechnologien -- Übungsblatt 1}}
\end{center}
\section*{Aufgabe 6: Mietwagenformular}
\textbf{Aufgabe:} Das Formular in \texttt{ex6.html} um Kontakt- und
Mietdaten sowie Fahrzeugwünsche ergänzen. HTML-Prüfungen verwenden,
ohne CSS oder JavaScript.

\medskip
\textbf{Ergebnis:}
\begin{itemize}
  \item Anrede, Name, E-Mail und Telefonnummer.
  \item Abhol- und Rückgabedatum, Abholort und Passagierzahl.
  \item Farbauswahl und sechs Fahrzeugtypen als Radiobuttons.
  \item Optionaler Kommentar mit maximal 512 Zeichen.
  \item Labels, Feldgruppen und Beispieltexte für Textfelder.
  \item Pflichtfelder, E-Mail-Prüfung, Telefonmuster und positive ganze Passagierzahl.
  \item Farbfeld mit Vorgabewert; damit wird immer ein Farbwert übermittelt.
  \item Absenden über \texttt{mailto} mit benannten Feldern und Reset-Button.
\end{itemize}
\textbf{Test:} Über Apache unter
\url{http://localhost/webtech-aufgabe6/ex6.html} erreichbar,
Status \texttt{200 OK}. Eingabeprüfungen und Zurücksetzen im Browser geprüft.

\medskip
Der E-Mail-Versand benötigt einen eingerichteten Mailclient und wurde
nicht ausgeführt. Die Reihenfolge von Abhol- und Rückgabedatum wird
ohne JavaScript nicht dynamisch geprüft.
\medskip
{\small \textbf{Hilfsmittel:} Codex für Erstellung, Prüfung und Dokumentation.}
\end{document}
